% !TeX root = ../presentacion.tex
% ===========================================================================
%  1. El encargo — 0:45
%
%  Defiende «Alcance funcional implementado» (25 %), el criterio de mayor
%  peso. Aquí solo se enuncia lo prometido; que funciona se demuestra en la
%  demo, no se afirma en una lámina.
%
%  Una sola lámina a propósito: el docente escribió el documento de alcance,
%  no hay que explicárselo. Hay que demostrar que se leyó.
% ===========================================================================

\bloque{El encargo}{0:45}{\exponeArquitectura}%
       {Alcance funcional implementado (25\,\%)}

\begin{frame}{Lo que pedía el documento de alcance}
  \begin{columns}[T]
    \begin{column}{0.56\textwidth}
      Un club con canchas de \textbf{pádel, tenis y básquet}. Dos roles:
      \begin{itemize}\footnotesize
        \item \textbf{Socio:} consulta disponibilidad, reserva un bloque y
              cancela lo suyo.
        \item \textbf{Administrador:} gestiona el catálogo y los bloqueos,
              cancela cualquier reserva, activa e inactiva usuarios y mide el
              uso real.
      \end{itemize}

      \vspace{4pt}
      {\footnotesize Sobre eso, \textbf{ocho reglas de negocio}
      (\rn{01}--\rn{08}) y \textbf{seis criterios de aceptación} (§7).
      Todo lo que sigue está atado a ese documento.}
    \end{column}
    \begin{column}{0.40\textwidth}
      \begin{block}{Fuera de alcance (§3.5)}
        \footnotesize
        Pagos en línea · notificaciones · reservas recurrentes · torneos ·
        app móvil nativa · BI avanzado.
      \end{block}
      {\scriptsize\color{textoSuave} No lo construimos, y es deliberado:
      está delimitado así en el propio documento.}
    \end{column}
  \end{columns}
\end{frame}
